\subsection{Suspended Execution and Runtime-Managed Control Flow}

The previous subsection showed that functions can become runtime values: they can be passed as callbacks, created anonymously, and preserved together with captured lexical state through closures. This subsection moves one step further. It examines what happens when execution itself can be paused, stored, and resumed later. At that point, the runtime is no longer only calling functions; it is managing unfinished computations.

In the traditional native C model, an ordinary function call creates a stack frame, runs, and eventually returns. Once the function returns, its stack frame is gone. Generators and coroutines break this simple lifecycle. A generator can execute part of its body, produce a value, suspend, and later continue from the same point. A coroutine can suspend not only to produce a value, but also to cooperate with other in-flight tasks. This requires the runtime to preserve execution state: local variables, instruction position, pending values, and sometimes a larger logical call context.

These mechanisms belong together because they all depend on the same architectural shift: control flow is no longer only a linear call-and-return chain on the native stack. Some execution state is moved into runtime-managed objects. In Python this may be a generator object, coroutine object, task, frame, or future. In JavaScript it may be a generator object, promise continuation, or async function state. In PHP it may be a generator or fiber. In C, where the language does not provide this directly, the programmer must emulate it with explicit state machines, structs, callbacks, or non-portable stack-switching mechanisms.

This is also where the distinction between concurrency and parallelism becomes necessary. A single thread can manage many suspended tasks by interleaving them: one task runs until it waits, then another task runs, then the first resumes later. This is concurrency without parallelism. It is especially useful for I/O-bound programs, where much time is spent waiting for files, sockets, timers, HTTP responses, user input, or database operations.

The purpose of this subsection is to connect generators, coroutines, and asynchronous event loops to the memory and execution model already established earlier. These are not magical language features. They are runtime control-flow mechanisms built on saved execution state. Understanding where the stack stops, where the runtime object begins, and who resumes the computation is the key to understanding lazy sequences, cooperative multitasking, and single-threaded asynchronous systems.


\subsubsection{Generators: Lazy Sequence Evaluation, In-Flight Stream Interfaces, and State-Preserving Suspended Subroutines}

A \textbf{generator} is a callable execution unit that can produce a sequence of values over time without computing and storing the whole sequence at once. The central idea is suspension. A normal function starts, runs until it returns or raises an exception, and then its execution frame disappears. A generator starts, runs until it reaches a suspension point, returns one value to the caller, and keeps enough internal state to continue later from the same point.

The basic model is:

\begin{codebox}[]{python}{Generators: Lazy Sequence Evaluation, In-Flight Stream Interfaces, and State-...}{sus22}
caller asks for next value
    generator frame resumes
        runs until yield
        produces value
        suspends
caller receives value

caller asks again
    same generator frame resumes after previous yield
        runs until next yield
        produces value
        suspends again
\end{codebox}


This is different from returning a list. A list is materialized: all elements are stored in memory before the caller consumes them. A generator is lazy: values are produced only when requested.

For example, this creates all values immediately:

\begin{verbatim}
values = [1, 2, 3, 4, 5]
\end{verbatim}

A generator-like object can produce them one at a time:

\begin{verbatim}
next -> 1
next -> 2
next -> 3
next -> 4
next -> 5
finished
\end{verbatim}

The main useful applications are:

\begin{itemize}
    \item processing large files line by line;
    \item producing infinite or very large sequences;
    \item building data-processing pipelines;
    \item avoiding unnecessary list materialization;
    \item representing streams of events, rows, records, tokens, or packets;
    \item separating value production from value consumption.
\end{itemize}

\paragraph{Python / CPython: generator functions and suspended Python frames.}

In Python, a function containing \texttt{yield} becomes a generator function. Calling it does not run the body immediately. It creates a generator object.

\begin{codebox}[]{python}{Python / CPython: generator functions and suspended Python frames.}{sus26}
def count_up_to(limit):
    current = 1

    while current <= limit:
        yield current
        current += 1


g = count_up_to(3)

print(next(g))  # 1
print(next(g))  # 2
print(next(g))  # 3
\end{codebox}


The execution sequence is:

\begin{codebox}[]{python}{Python / CPython: generator functions and suspended Python frames.}{sus25}
count_up_to(3)
    creates generator object
    body has not run yet

next(g)
    creates/resumes generator frame
    current = 1
    reaches yield current
    returns 1 and suspends

next(g)
    resumes after yield
    current += 1
    reaches yield current
    returns 2 and suspends
\end{codebox}


At the CPython level, the generator object owns a suspended execution state. The old Python frame is not simply destroyed after the first \texttt{yield}. Its local variables and bytecode position remain associated with the generator.

Conceptually:

\begin{verbatim}
g
|
v
generator object
+----------------------------------+
| code object                      |
| suspended bytecode position      |
| local variable: current          |
| local variable: limit            |
| evaluation stack state           |
+----------------------------------+
\end{verbatim}

When the generator is suspended, the native C stack used by CPython is not preserved as a raw stack segment for that generator. The generator's Python-level state is stored in interpreter-managed structures. Later, when \texttt{next(g)} is called again, CPython re-enters that saved Python execution state.

During active execution, the Python frame chain may look like:

\begin{verbatim}
top Python frame
+-----------------------------+
| frame: count_up_to          |
+-----------------------------+
| frame: caller/module body   |
+-----------------------------+
bottom
\end{verbatim}

After \texttt{yield}, the caller continues, and the generator frame is suspended inside the generator object:

\begin{verbatim}
active caller frame
+-----------------------------+
| frame: caller/module body   |
+-----------------------------+

suspended generator object
+-----------------------------+
| saved frame: count_up_to    |
+-----------------------------+
\end{verbatim}

This distinction is essential. A generator does not keep running in the background. It is paused. It advances only when the caller requests another value.

A practical Python example is reading a large file lazily:

\begin{codebox}[]{python}{Python / CPython: generator functions and suspended Python frames.}{sus24}
def non_empty_lines(path):
    with open(path, "r", encoding="utf-8") as file:
        for line in file:
            line = line.strip()

            if line:
                yield line


for line in non_empty_lines("log.txt"):
    print(line)
\end{codebox}


The useful application is memory-safe file processing. The file is not loaded entirely into a list. Each useful line is produced when the loop asks for it.

The \texttt{for} loop drives the generator through the iterator protocol:

\begin{verbatim}
iter(generator)
next(generator)
next(generator)
...
StopIteration
\end{verbatim}

When the generator finishes normally, Python raises \texttt{StopIteration} internally to signal exhaustion to the loop machinery.

\paragraph{JavaScript / V8: generator functions and resumable function state.}

JavaScript supports generator functions using \texttt{function*} and \texttt{yield}.

\begin{codebox}[]{javascript}{JavaScript / V8: generator functions and resumable function state.}{sus13}
function* countUpTo(limit) {
    let current = 1;

    while (current <= limit) {
        yield current;
        current += 1;
    }
}

const g = countUpTo(3);

console.log(g.next()); // { value: 1, done: false }
console.log(g.next()); // { value: 2, done: false }
console.log(g.next()); // { value: 3, done: false }
console.log(g.next()); // { value: undefined, done: true }
\end{codebox}


The useful application is lazy traversal, custom iterable objects, streaming transformations, and controlled production of values.

Calling the generator function does not execute its body immediately:

\begin{verbatim}
countUpTo(3)
    returns generator object
    body not executed yet
\end{verbatim}

Each \texttt{next()} resumes the generator until the next \texttt{yield}.

The V8-level conceptual structure is:

\begin{verbatim}
generator object
+----------------------------------+
| generator function code          |
| suspended instruction position   |
| local binding: current           |
| local binding: limit             |
| generator state                  |
+----------------------------------+
\end{verbatim}

When the generator yields, its JavaScript call frame is not kept as an ordinary active stack frame. The resumable state is stored by the engine in heap-managed generator structures. Later, \texttt{next()} reconstructs or resumes the required execution state.

A practical JavaScript example is lazy ID generation:

\begin{codebox}[]{javascript}{JavaScript / V8: generator functions and resumable function state.}{sus12}
function* idGenerator(prefix) {
    let id = 1;

    while (true) {
        yield `${prefix}-${id}`;
        id += 1;
    }
}

const ids = idGenerator("user");

console.log(ids.next().value); // user-1
console.log(ids.next().value); // user-2
console.log(ids.next().value); // user-3
\end{codebox}


This is useful because the sequence is infinite. A list of all possible IDs cannot be materialized. The generator stores only the current state.

The active stack during one \texttt{next()} call is:

\begin{verbatim}
top JS frame
+-----------------------------+
| frame: idGenerator          |
+-----------------------------+
| frame: caller               |
+-----------------------------+
bottom
\end{verbatim}

After \texttt{yield}, the generator frame is suspended and removed from the ordinary active call stack:

\begin{verbatim}
caller continues

generator object remembers:
    prefix
    id
    suspended position
\end{verbatim}

\paragraph{PHP / Zend Engine: generators with \texttt{yield}.}

PHP also supports generators with \texttt{yield}. A function containing \texttt{yield} returns a generator object instead of executing to completion immediately.

\begin{codebox}[]{php}{PHP / Zend Engine: generators with yield.}{sus21}
<?php

function count_up_to($limit) {
    $current = 1;

    while ($current <= $limit) {
        yield $current;
        $current += 1;
    }
}

foreach (count_up_to(3) as $value) {
    echo $value . PHP_EOL;
}
\end{codebox}


The useful application is streaming data without building large arrays.

For example, reading a file lazily:

\begin{codebox}[]{php}{PHP / Zend Engine: generators with yield.}{sus20}
<?php

function non_empty_lines($path) {
    $file = fopen($path, "r");

    while (($line = fgets($file)) !== false) {
        $line = trim($line);

        if ($line !== "") {
            yield $line;
        }
    }

    fclose($file);
}

foreach (non_empty_lines("log.txt") as $line) {
    echo $line . PHP_EOL;
}
\end{codebox}


Without a generator, one might build an array of all lines and return it. With a generator, each line is produced only when the \texttt{foreach} loop requests it.

At the Zend Engine level, the generator object preserves the function's execution state: local variables, current opcode position, and yielded value state. The ordinary active call frame does not remain on the native C stack while the generator is suspended. The engine stores the resumable state in runtime-managed generator structures.

Conceptually:

\begin{verbatim}
Generator object
+----------------------------------+
| compiled opcodes                 |
| suspended opcode position        |
| local variable: $current         |
| local variable: $limit           |
| current yielded value            |
+----------------------------------+
\end{verbatim}

During one resume, the PHP frame is active:

\begin{verbatim}
top PHP frame
+-----------------------------+
| frame: count_up_to          |
+-----------------------------+
| frame: foreach/script body  |
+-----------------------------+
bottom
\end{verbatim}

After \texttt{yield}, the script body continues and the generator frame is suspended inside the generator object.

\paragraph{Java: no native generator syntax, but iterator objects emulate lazy production.}

Java does not have Python-style generator functions in the core language. The usual Java approach is to represent lazy production with an \texttt{Iterator}, \texttt{Iterable}, \texttt{Stream}, or a custom object that stores iteration state explicitly.

A manual iterator example:

\begin{codebox}[]{java}{Java: no native generator syntax, but iterator objects emulate lazy production.}{sus08}
import java.util.Iterator;
import java.util.NoSuchElementException;

class CountUpTo implements Iterator<Integer> {
    private int current = 1;
    private final int limit;

    CountUpTo(int limit) {
        this.limit = limit;
    }

    public boolean hasNext() {
        return current <= limit;
    }

    public Integer next() {
        if (!hasNext()) {
            throw new NoSuchElementException();
        }

        int value = current;
        current += 1;
        return value;
    }
}

public class Main {
    public static void main(String[] args) {
        Iterator<Integer> it = new CountUpTo(3);

        while (it.hasNext()) {
            System.out.println(it.next());
        }
    }
}
\end{codebox}


The useful application is still lazy traversal. The values are not stored in a list. The object stores the state needed to compute the next value.

However, the implementation is different from Python generators. Java does not suspend a method frame at a \texttt{yield} statement here. The state is manually stored in object fields:

\begin{verbatim}
CountUpTo object
+-----------------------------+
| current                     |
| limit                       |
+-----------------------------+
\end{verbatim}

Each call to \texttt{next()} creates an ordinary JVM method frame:

\begin{verbatim}
top JVM frame
+-----------------------------+
| frame: CountUpTo.next       |
+-----------------------------+
| frame: Main.main            |
+-----------------------------+
bottom
\end{verbatim}

When \texttt{next()} returns, that frame disappears. The persistent state is not the old stack frame. It is the object field \texttt{current} stored on the heap.

Java streams can also express lazy pipelines:

\begin{codebox}[]{java}{Java: no native generator syntax, but iterator objects emulate lazy production.}{sus07}
import java.util.stream.IntStream;

public class Main {
    public static void main(String[] args) {
        IntStream.iterate(1, x -> x + 1)
                 .limit(3)
                 .forEach(System.out::println);
    }
}
\end{codebox}


The useful application is declarative lazy data processing. The stream pipeline stores transformation logic and pulls values as needed. Internally, Java stream implementations use objects, spliterators, lambdas, and library machinery, not preserved Java method stack frames in the Python generator sense.

\paragraph{C: emulating generators with explicit state machines.}

Standard C has no generator syntax and no automatic suspension of function frames. A C function call normally runs until it returns. Once it returns, its stack frame is gone. To emulate a generator, the programmer must store state explicitly.

A simple C generator-like object:

\begin{codebox}[]{c}{C: emulating generators with explicit state machines.}{sus03}
#include <stdio.h>

struct counter {
    int current;
    int limit;
};

int counter_next(struct counter *c, int *out) {
    if (c->current > c->limit) {
        return 0;
    }

    *out = c->current;
    c->current += 1;
    return 1;
}

int main(void) {
    struct counter c = { .current = 1, .limit = 3 };

    int value;
    while (counter_next(&c, &value)) {
        printf("%d\n", value);
    }

    return 0;
}
\end{codebox}


The useful application is lazy production in embedded systems, parsers, stream processors, or low-level libraries where allocating a full list would be wasteful.

The state is explicit:

\begin{verbatim}
struct counter
+-----------------------------+
| current                     |
| limit                       |
+-----------------------------+
\end{verbatim}

The stack behavior is ordinary C behavior:

\begin{verbatim}
main()
    calls counter_next()
        counter_next frame is created
        computes one value
        returns
    counter_next frame disappears

main()
    calls counter_next() again
        new counter_next frame is created
        uses saved state from struct
        returns
\end{verbatim}

There is no suspended C stack frame. The persistent state lives in the \texttt{struct counter}, not in the old call frame.

For more complex control flow, C programmers sometimes build manual state machines with \texttt{switch}:

\begin{codebox}[]{c}{C: emulating generators with explicit state machines.}{sus02}
#include <stdio.h>

struct generator {
    int state;
    int i;
};

int gen_next(struct generator *g, int *out) {
    switch (g->state) {
        case 0:
            g->i = 1;

        case 1:
            if (g->i <= 3) {
                *out = g->i;
                g->i += 1;
                g->state = 1;
                return 1;
            }

            g->state = 2;

        case 2:
            return 0;
    }

    return 0;
}
\end{codebox}


This uses a stored program counter-like state variable. It is a crude manual version of what a generator runtime does automatically: remember where execution should continue.

Duff's Device is related only as a demonstration that C \texttt{switch} and loop control can be arranged in unusual ways to resume at different points in a block. It is not a true generator feature. However, the same idea of storing a state label and jumping into the right part of a function appears in coroutine/generator emulation macros.

A famous style of C coroutine macro uses \texttt{switch} and \texttt{\_\_LINE\_\_} to remember the resume point. Conceptually:

\begin{verbatim}
state says where to resume
switch(state):
    case previous_yield_location:
        continue execution there
\end{verbatim}

This is powerful but fragile. It depends on source layout tricks and does not preserve ordinary C local variables automatically unless they are stored outside the function frame.

\paragraph{Bash: stream-like generation through output and pipelines.}

Bash has no generator objects or suspended shell function frames in the Python sense. However, shell programming often uses stream-like lazy processing through pipes. A command produces output gradually, and another command consumes it gradually.

Example:

\begin{verbatim}
generate_numbers() {
    i=1

    while [ "$i" -le 3 ]; do
        echo "$i"
        i=$((i + 1))
    done
}

generate_numbers | while read -r value; do
    echo "value: $value"
done
\end{verbatim}

The useful application is pipeline processing. The producer does not need to store all output in an array. It writes lines. The consumer reads lines.

The implementation is not a generator frame. In many shells, the two sides of a pipeline run in separate processes or subshells:

\begin{verbatim}
process 1:
    generate_numbers writes to pipe

process 2:
    while read consumes from pipe
\end{verbatim}

The operating system pipe provides buffering. The producer and consumer are coordinated by I/O blocking, not by a language-level \texttt{yield} instruction. If the pipe buffer fills, the producer blocks. If the pipe is empty, the consumer blocks.

A Bash function can also be called repeatedly, but it does not preserve a suspended frame automatically. Persistent state usually lives in variables, files, or external processes, not in a generator object.

\paragraph{Implementation comparison.}

Generators and generator-like mechanisms all solve the same practical problem: produce values gradually while remembering progress. The implementation differs sharply by language.

\begin{center}
\begin{tabular}{|p{0.15\textwidth}|p{0.34\textwidth}|p{0.35\textwidth}|}
\hline
\textbf{Language} & \textbf{Generator form} & \textbf{State and stack behavior} \\
\hline
Python / CPython & Function with \texttt{yield} & Generator object preserves Python frame state, locals, and bytecode position; native C stack is not preserved while suspended \\
\hline
JavaScript / V8 & \texttt{function*} with \texttt{yield} & Generator object preserves engine-managed continuation state and lexical environment; JS stack frame is suspended into heap-managed state \\
\hline
PHP / Zend & Function with \texttt{yield} & Generator object preserves opcode position and local variables; ordinary active call frame is not kept on native stack \\
\hline
Java & \texttt{Iterator}, \texttt{Iterable}, \texttt{Stream} & No built-in suspended method frame; state stored manually or by library objects on heap \\
\hline
C & Manual state machine, struct, function pointer patterns & No suspended C frame; programmer stores state explicitly in structs, often heap-allocated if it must escape \\
\hline
Bash & Pipes, command output streams, repeated reads & No generator object; OS pipes and child processes provide streaming behavior \\
\hline
\end{tabular}
\end{center}

The most important stack rule is this:

\begin{quote}
A real generator does not usually keep the old native call stack alive. It stores the resumable language-level state in a heap-managed generator object or equivalent runtime structure.
\end{quote}

That is why a generator can survive after the function call that created it appears to have returned. The call did not return a final result. It returned a resumable object.

The difference between a generator and a list is therefore not only memory size. It is execution architecture. A list is stored data. A generator is suspended computation.

\begin{verbatim}
list:
    all values already materialized

generator:
    code + saved state + ability to produce next value
\end{verbatim}

This prepares the transition to coroutines. A generator produces values outward to a caller. A coroutine generalizes the idea: execution can suspend and resume as part of a larger cooperative control-flow system, sometimes receiving values, waiting for events, or interleaving multiple in-flight tasks.



\subsubsection{Coroutines: Non-Preemptive Cooperative Tasking, Symmetric Yield Transfers, and Context Interleaving}

A \textbf{coroutine} is a resumable execution unit that can suspend itself and later continue from the same logical point. A generator is already a restricted coroutine: it yields values outward to its caller. A fuller coroutine can also receive values, cooperate with other execution units, and participate in a scheduler-controlled flow where several in-progress computations take turns.

The key word is \textbf{cooperative}. A coroutine is not usually interrupted by the operating system in the middle of an arbitrary instruction. It runs until it reaches an explicit suspension point such as \texttt{yield}, \texttt{await}, \texttt{send}, or a runtime-specific operation. At that point, it gives control back to a caller or scheduler.

\begin{verbatim}
coroutine A runs
    reaches suspension point
scheduler resumes coroutine B
    coroutine B reaches suspension point
scheduler resumes coroutine A
    coroutine A continues from where it stopped
\end{verbatim}

This is non-preemptive tasking. The coroutine must cooperate. If it enters a long CPU loop and never suspends, other coroutines in the same runtime scheduler do not progress.

\paragraph{Python: generator-based coroutines with \texttt{send}.}

Before modern \texttt{async}/\texttt{await}, Python generators could already behave like simple coroutines because values can be sent back into them.

\begin{codebox}[]{python}{Python: generator-based coroutines with send.}{sus30}
def accumulator():
    total = 0

    while True:
        value = yield total
        total += value


acc = accumulator()

print(next(acc))      # starts coroutine, yields 0
print(acc.send(5))    # sends 5, yields 5
print(acc.send(7))    # sends 7, yields 12
\end{codebox}


The useful application is stateful stream processing. The coroutine keeps \texttt{total} between resumptions without using a global variable.

At the CPython level, \texttt{accumulator()} creates a generator/coroutine-like object. The local variable \texttt{total}, the bytecode position, and the suspended execution state are stored inside interpreter-managed structures.

\begin{verbatim}
acc
|
v
generator/coroutine object
+----------------------------------+
| code object                      |
| suspended bytecode position      |
| local variable: total            |
| slot receiving sent value        |
+----------------------------------+
\end{verbatim}

The native C stack is not kept alive while the coroutine is suspended. When \texttt{yield} executes, control returns to the caller. The Python frame state is preserved by the runtime object. When \texttt{send()} is called, CPython resumes that saved frame and injects the sent value as the result of the suspended \texttt{yield} expression.

During active execution:

\begin{verbatim}
top Python frame
+-----------------------------+
| frame: accumulator          |
+-----------------------------+
| frame: caller/module body   |
+-----------------------------+
bottom
\end{verbatim}

After suspension:

\begin{verbatim}
active caller frame remains

acc object stores:
    total
    bytecode position
    suspended yield point
\end{verbatim}

A small cooperative scheduler can be sketched with generators:

\begin{codebox}[]{python}{Python: generator-based coroutines with send.}{sus29}
def task(name, count):
    for i in range(count):
        print(name, i)
        yield


tasks = [task("A", 3), task("B", 3)]

while tasks:
    current = tasks.pop(0)

    try:
        next(current)
        tasks.append(current)
    except StopIteration:
        pass
\end{codebox}


The output interleaves the tasks:

\begin{verbatim}
A 0
B 0
A 1
B 1
A 2
B 2
\end{verbatim}

This is concurrency by interleaving, not parallelism. One Python task runs at a time. The scheduler advances each task until its next \texttt{yield}.

\paragraph{JavaScript: generator coroutines and explicit resumption.}

JavaScript generators can also receive values through \texttt{next(value)}.

\begin{codebox}[]{javascript}{JavaScript: generator coroutines and explicit resumption.}{sus10}
function* accumulator() {
    let total = 0;

    while (true) {
        const value = yield total;
        total += value;
    }
}

const acc = accumulator();

console.log(acc.next().value);   // 0
console.log(acc.next(5).value);  // 5
console.log(acc.next(7).value);  // 12
\end{codebox}


The useful application is resumable state machines, parser pipelines, controlled data producers, and cooperative simulations.

V8 stores the generator's resumable state in a heap-managed generator object:

\begin{verbatim}
generator object
+----------------------------------+
| generator function code          |
| suspended instruction position   |
| local binding: total             |
| generator state                  |
+----------------------------------+
\end{verbatim}

When the generator is suspended, the ordinary JavaScript call stack is not still active. A later \texttt{next(...)} call creates/resumes the engine frame needed to continue from the saved point.

JavaScript's common asynchronous programming model later developed around promises and \texttt{async}/\texttt{await}, but generator coroutines are the simpler mechanical model: resume explicitly, run until \texttt{yield}, store state, return control.

\paragraph{PHP: generator coroutines with \texttt{send}.}

PHP generators also support sending values back into a suspended generator.

\begin{codebox}[]{php}{PHP: generator coroutines with send.}{sus17}
<?php

function accumulator() {
    $total = 0;

    while (true) {
        $value = yield $total;
        $total += $value;
    }
}

$acc = accumulator();

echo $acc->current() . PHP_EOL; // starts, yields 0

$acc->send(5);
echo $acc->current() . PHP_EOL; // 5

$acc->send(7);
echo $acc->current() . PHP_EOL; // 12
\end{codebox}


The useful application is cooperative pipelines and older generator-based async frameworks. Before native async syntax was common in many ecosystems, generator suspension was often used as a building block for user-space schedulers.

At the Zend Engine level, the generator object stores the suspended opcode position and local variables. The old PHP call frame is not kept as an active native stack frame.

\begin{verbatim}
Generator object
+----------------------------------+
| compiled opcodes                 |
| suspended opcode position        |
| local variable: $total           |
| current yielded value            |
+----------------------------------+
\end{verbatim}

Each resume re-enters the generator's saved execution state until it yields again or finishes.

\paragraph{Lua: classic coroutine-style cooperative transfer.}

Lua is worth mentioning because it has direct coroutine support and is often used as the clean textbook example of cooperative coroutines.

\begin{codebox}[]{python}{Lua: classic coroutine-style cooperative transfer.}{sus23}
co = coroutine.create(function ()
    for i = 1, 3 do
        print("coroutine", i)
        coroutine.yield()
    end
end)

coroutine.resume(co)
print("main")
coroutine.resume(co)
print("main again")
coroutine.resume(co)
\end{codebox}


The useful application is game scripting, embedded scripting, cooperative task systems, and lightweight state machines.

Lua coroutines are stackful at the Lua-runtime level: each coroutine has its own Lua call stack managed by the Lua VM. When it yields, that Lua stack is preserved by the runtime. This is stronger than Python-style generator suspension, which is usually a single suspended generator frame rather than a whole independent call stack of arbitrary depth.

Conceptually:

\begin{verbatim}
Lua coroutine object
+----------------------------------+
| Lua call stack                   |
| current instruction position     |
| local variables                  |
| coroutine status                 |
+----------------------------------+
\end{verbatim}

The operating system still sees one native process and usually one native thread unless the host program creates more. Lua coroutine switching is runtime-level switching, not OS preemption.

\paragraph{Java: no direct language coroutine, but runtime-managed alternatives.}

Java does not have a direct \texttt{yield}-based coroutine syntax equivalent to Python generators or Lua coroutines. Similar behavior is usually expressed through objects, iterators, streams, callbacks, futures, or more recently virtual threads.

A manual cooperative task can be represented as an object with explicit state:

\begin{verbatim}
interface StepTask {
    boolean step();
}

class CounterTask implements StepTask {
    private final String name;
    private int current = 0;
    private final int limit;

    CounterTask(String name, int limit) {
        this.name = name;
        this.limit = limit;
    }

    public boolean step() {
        if (current >= limit) {
            return false;
        }

        System.out.println(name + " " + current);
        current++;
        return true;
    }
}
\end{verbatim}

A scheduler can interleave tasks:

\begin{codebox}[]{java}{Java: no direct language coroutine, but runtime-managed alternatives.}{sus06}
import java.util.*;

public class Main {
    public static void main(String[] args) {
        Queue<StepTask> tasks = new ArrayDeque<>();

        tasks.add(new CounterTask("A", 3));
        tasks.add(new CounterTask("B", 3));

        while (!tasks.isEmpty()) {
            StepTask task = tasks.remove();

            if (task.step()) {
                tasks.add(task);
            }
        }
    }
}
\end{codebox}


The useful application is simulations, event loops, and manually scheduled state machines.

There is no suspended Java method frame here. Each call to \texttt{step()} creates an ordinary JVM frame, then returns. Persistent state lives in object fields:

\begin{verbatim}
CounterTask object
+-----------------------------+
| name                        |
| current                     |
| limit                       |
+-----------------------------+
\end{verbatim}

Java virtual threads are different. They are JVM-managed lightweight threads that can be parked and resumed by the runtime. They are useful for writing blocking-looking server code with much lower thread overhead than traditional platform threads. However, they are not the same as lexical generator coroutines. They are scheduled by the JVM and integrate with blocking operations, while source code still looks like ordinary method calls.

\paragraph{C: coroutine emulation with explicit state machines or separate stacks.}

Standard C has no coroutine feature. A C function call uses the native call stack and normally runs until it returns. To emulate coroutines portably, the programmer stores the resume state manually.

A cooperative task as an explicit state machine:

\begin{codebox}[]{c}{C: coroutine emulation with explicit state machines or separate stacks.}{sus01}
#include <stdio.h>

struct task {
    const char *name;
    int state;
    int current;
    int limit;
};

int task_step(struct task *t) {
    switch (t->state) {
        case 0:
            t->current = 0;
            t->state = 1;

        case 1:
            if (t->current < t->limit) {
                printf("%s %d\n", t->name, t->current);
                t->current++;
                return 1;
            }

            t->state = 2;

        case 2:
            return 0;
    }

    return 0;
}

int main(void) {
    struct task a = {"A", 0, 0, 3};
    struct task b = {"B", 0, 0, 3};

    while (task_step(&a) | task_step(&b)) {
        /* keep stepping while either task can run */
    }

    return 0;
}
\end{codebox}


The useful application is embedded systems, protocol parsers, cooperative schedulers, and low-level event loops.

The stack behavior is ordinary C:

\begin{verbatim}
main()
    calls task_step(&a)
        task_step frame created
        one step executes
        task_step returns
    frame disappears

main()
    calls task_step(&b)
        new task_step frame created
        one step executes
        task_step returns
    frame disappears
\end{verbatim}

No C stack frame is suspended. The resume point is stored in \texttt{state}; local state is stored in the \texttt{struct task}.

Duff's Device is relevant as a C control-flow trick because it shows how \texttt{switch} labels can be interleaved with loop structure. Coroutine macros in C often use a similar idea: store a resume label, then enter a \texttt{switch} at the right case on the next call. But this is not true native coroutine support. It is manual transformation of a coroutine into a state machine.

Non-portable C coroutine systems can use separate stacks, assembly, POSIX \texttt{ucontext} APIs, Windows fibers, or library-level context switching. In that case, each coroutine may have its own native stack region:

\begin{verbatim}
coroutine A
+-----------------------------+
| native stack for A          |
+-----------------------------+

coroutine B
+-----------------------------+
| native stack for B          |
+-----------------------------+
\end{verbatim}

Switching saves CPU registers and stack pointers, then restores another coroutine's saved context. This is much closer to stackful coroutines, but it is outside portable standard C.

\paragraph{Bash: cooperative interleaving through processes, pipes, and jobs.}

Bash has no coroutine object or suspended shell frame. It can still interleave work through loops, background jobs, pipes, and process coordination.

A simple cooperative-looking loop:

\begin{verbatim}
task_a_step() {
    echo "A $1"
}

task_b_step() {
    echo "B $1"
}

for i in 0 1 2; do
    task_a_step "$i"
    task_b_step "$i"
done
\end{verbatim}

This is manual interleaving, not coroutine suspension. Each shell function call runs and returns normally. Its frame disappears.

Bash can also run background jobs:

\begin{verbatim}
long_task_a &
pid_a=$!

long_task_b &
pid_b=$!

wait "$pid_a"
wait "$pid_b"
\end{verbatim}

This is not runtime-level coroutine scheduling. Bash asks the operating system to run child processes. Each child has its own process state and native stack. The OS scheduler, not Bash coroutine logic, interleaves them.

Pipelines also interleave processes:

\begin{verbatim}
producer | consumer
\end{verbatim}

The producer and consumer may run concurrently as separate processes connected by an OS pipe. This is useful for streaming, but it is process-level coordination, not a Bash coroutine frame.

\paragraph{Symmetric and asymmetric transfer.}

Many common generators are \textbf{asymmetric}: they yield back to their immediate caller or scheduler. The caller resumes them later.

\begin{verbatim}
caller -> generator
generator -> caller
caller -> generator
generator -> caller
\end{verbatim}

A more symmetric coroutine model allows coroutine A to transfer control directly to coroutine B, or to a scheduler that chooses the next coroutine.

\begin{verbatim}
coroutine A -> coroutine B
coroutine B -> coroutine C
coroutine C -> coroutine A
\end{verbatim}

Many practical systems implement this through a scheduler:

\begin{verbatim}
coroutine A -> scheduler
scheduler -> coroutine B
coroutine B -> scheduler
scheduler -> coroutine C
\end{verbatim}

The scheduler keeps a queue of resumable tasks. Each task runs until it suspends. No task is preempted in the middle of ordinary computation.

\paragraph{Comparison.}

\begin{center}
\begin{tabular}{|p{0.15\textwidth}|p{0.34\textwidth}|p{0.35\textwidth}|}
\hline
\textbf{Language} & \textbf{Coroutine-like mechanism} & \textbf{Stack and state behavior} \\
\hline
Python / CPython & Generators with \texttt{send}; later \texttt{async}/\texttt{await} & Runtime stores suspended Python frame state; native C stack is not preserved while suspended \\
\hline
JavaScript / V8 & Generators with \texttt{next(value)}; async functions later & V8 stores generator/coroutine state in heap-managed engine structures \\
\hline
PHP / Zend & Generators with \texttt{send} & Zend generator object stores opcode position and locals \\
\hline
Lua & Built-in coroutines & Lua VM preserves coroutine call stack and state inside coroutine object \\
\hline
Java & Manual state objects, futures, virtual threads & Ordinary methods do not suspend; state is stored in objects, or JVM manages virtual-thread continuation state \\
\hline
C & Manual state machines; nonportable fibers/context switching & Portable C stores state explicitly; stackful libraries use separate stacks and saved register contexts \\
\hline
Bash & Manual interleaving, jobs, pipelines & No coroutine frames; shell functions return normally, or OS processes are scheduled separately \\
\hline
\end{tabular}
\end{center}

The important architectural point is that a coroutine separates \emph{returning control} from \emph{finishing execution}. A normal function returns and is done. A coroutine yields, pauses, and can later continue. This is possible only if the runtime or programmer stores enough state: current instruction position, local variables, pending values, and sometimes an entire call stack.

This prepares the next distinction: coroutine interleaving can create concurrency without parallelism. Several tasks may be in progress at once from the program's point of view, even though only one task is executing on the CPU at any exact moment.



\subsubsection{Concurrency Without Parallelism: Interleaving Waiting Tasks on One Thread}

\textbf{Concurrency} means that several tasks are in progress during the same time interval. \textbf{Parallelism} means that several tasks are executing at the exact same time on different CPU cores or hardware threads. These are not the same idea.

A single-threaded event loop can be concurrent without being parallel. It can manage many tasks that are waiting for timers, network data, file descriptors, user input, or other external events. At any exact moment, only one piece of user code runs on the thread. However, while one task is waiting, the event loop can let another task advance.

The general shape is:

\begin{verbatim}
task A starts
task A reaches waiting point
task A is suspended

task B starts
task B reaches waiting point
task B is suspended

external event for task A becomes ready
task A resumes

external event for task B becomes ready
task B resumes
\end{verbatim}

The key is that waiting must be non-blocking. If a task blocks the only thread with a normal blocking operation, no other task can run.

\paragraph{Python / CPython: \texttt{asyncio} tasks on one event-loop thread.}

Python's \texttt{asyncio} model is a direct example of concurrency without parallelism. Multiple coroutines can be active, but only one coroutine executes Python bytecode at a time on the event-loop thread.

\begin{codebox}[]{python}{Python / CPython: asyncio tasks on one event-loop thread.}{sus28}
import asyncio


async def worker(name, delay):
    print(name, "start")

    await asyncio.sleep(delay)

    print(name, "done")


async def main():
    await asyncio.gather(
        worker("A", 1),
        worker("B", 1),
    )


asyncio.run(main())
\end{codebox}


The useful application is I/O-heavy programming: network clients, web servers, websocket handlers, database calls, crawlers, API clients, and other tasks where much time is spent waiting.

The output is approximately:

\begin{verbatim}
A start
B start
A done
B done
\end{verbatim}

Both tasks finish after about one second, not because they ran in parallel on two cores, but because \texttt{asyncio.sleep()} does not block the thread. It registers a timer and gives control back to the event loop.

The stack behavior is:

\begin{verbatim}
event loop resumes worker A
    worker A runs until await
    worker A frame is suspended in a Task object

event loop resumes worker B
    worker B runs until await
    worker B frame is suspended in a Task object

timer for A becomes ready
    event loop resumes worker A

timer for B becomes ready
    event loop resumes worker B
\end{verbatim}

At suspension time, the coroutine's Python frame is not active on the call stack. Its execution state is preserved in runtime-managed coroutine/task structures:

\begin{verbatim}
Task A
+----------------------------------+
| coroutine object                 |
| suspended frame                  |
| bytecode position at await       |
| local variables                  |
| awaited Future / timer           |
+----------------------------------+
\end{verbatim}

The wrong version is:

\begin{verbatim}
import time
import asyncio


async def bad_worker():
    time.sleep(5)      # blocks the whole event-loop thread
    print("done")
\end{verbatim}

Here, \texttt{time.sleep()} blocks the OS thread. No other coroutine can advance during that time. The correct async version is:

\begin{verbatim}
await asyncio.sleep(5)
\end{verbatim}

That suspends only the coroutine, not the whole event-loop thread.

\paragraph{JavaScript / V8: browser and Node.js event loops.}

JavaScript in browsers and Node.js is strongly associated with single-threaded concurrency. The engine runs one JavaScript call stack at a time, while the host environment manages timers, network events, DOM events, file events, and promise continuations.

Example with timers:

\begin{codebox}[]{javascript}{JavaScript / V8: browser and Node.js event loops.}{sus11}
console.log("start");

setTimeout(() => {
    console.log("timer A");
}, 1000);

setTimeout(() => {
    console.log("timer B");
}, 1000);

console.log("after scheduling");
\end{codebox}


The useful application is browser UI programming, AJAX/fetch completion handling, timers, Node.js servers, websocket servers, and event-driven desktop applications built on Electron-style runtimes.

The execution order is:

\begin{verbatim}
start
after scheduling
timer A
timer B
\end{verbatim}

The first stack turn registers the callbacks:

\begin{verbatim}
top-level script
    calls setTimeout for A
    host stores callback A
    calls setTimeout for B
    host stores callback B
top-level script finishes
\end{verbatim}

Later, when timers expire, the host event loop calls the callbacks:

\begin{verbatim}
event loop turn:
+-----------------------------+
| frame: callback A           |
+-----------------------------+
| frame: host dispatch        |
+-----------------------------+

next event loop turn:
+-----------------------------+
| frame: callback B           |
+-----------------------------+
| frame: host dispatch        |
+-----------------------------+
\end{verbatim}

The old registration stack is gone. The callbacks survive because the host/runtime keeps references to function objects.

AJAX/fetch has the same architecture:

\begin{verbatim}
fetch("/data.json")
    .then(response => response.json())
    .then(data => {
        console.log(data);
    });
\end{verbatim}

The network operation is not performed by continuously running JavaScript code. The host starts the operation, JavaScript returns to the event loop, and the callback resumes later when the result is ready.

\paragraph{PHP / Zend Engine: Fibers and user-space async frameworks.}

Modern PHP has \texttt{Fiber}, which allows cooperative suspension and resumption. A fiber is not an OS thread. It is a runtime-managed execution context.

\begin{codebox}[]{php}{PHP / Zend Engine: Fibers and user-space async frameworks.}{sus19}
<?php

$fiberA = new Fiber(function () {
    echo "A start\n";
    Fiber::suspend();
    echo "A done\n";
});

$fiberB = new Fiber(function () {
    echo "B start\n";
    Fiber::suspend();
    echo "B done\n";
});

$fiberA->start();
$fiberB->start();

$fiberA->resume();
$fiberB->resume();
\end{codebox}


\begin{iobox}{Expected Output}
A start \\
B start \\
A done \\
B done \\
\end{iobox}


The useful application is user-space async frameworks, cooperative I/O libraries, request pipelines, and systems that want to write code in a sequential-looking style while still yielding control at waiting points.

At the Zend Engine level, the fiber preserves PHP execution state. When \texttt{Fiber::suspend()} is called, the currently running fiber stops and returns control to the caller/scheduler. Its PHP-level call stack and local variables are preserved by the runtime.

Conceptually:

\begin{verbatim}
Fiber A
+----------------------------------+
| suspended PHP execution context  |
| local variables                  |
| current opcode position          |
+----------------------------------+

Fiber B
+----------------------------------+
| suspended PHP execution context  |
| local variables                  |
| current opcode position          |
+----------------------------------+
\end{verbatim}

Only one fiber runs at a time on the current thread. This is concurrency by cooperative interleaving, not CPU parallelism.

\paragraph{Java: single-thread event loops and callback dispatch.}

Java does not need coroutines to demonstrate concurrency without parallelism. A single-threaded event loop can run callbacks one at a time. GUI systems and server frameworks often use this idea. For example, Swing has an Event Dispatch Thread; frameworks such as Netty and Vert.x use event-loop-style execution.

A small simplified Java event loop:

\begin{codebox}[]{java}{Java: single-thread event loops and callback dispatch.}{sus09}
import java.util.ArrayDeque;
import java.util.Queue;

public class EventLoop {
    private final Queue<Runnable> queue = new ArrayDeque<>();

    public void callSoon(Runnable task) {
        queue.add(task);
    }

    public void run() {
        while (!queue.isEmpty()) {
            Runnable task = queue.remove();
            task.run();
        }
    }

    public static void main(String[] args) {
        EventLoop loop = new EventLoop();

        loop.callSoon(() -> System.out.println("A"));
        loop.callSoon(() -> System.out.println("B"));
        loop.callSoon(() -> System.out.println("C"));

        loop.run();
    }
}
\end{codebox}


The useful application is event dispatch: GUI actions, network events, timers, and framework callbacks.

The stack during one callback is:

\begin{verbatim}
top JVM frame
+-----------------------------+
| frame: lambda body          |
+-----------------------------+
| frame: Runnable.run         |
+-----------------------------+
| frame: EventLoop.run        |
+-----------------------------+
| frame: main                 |
+-----------------------------+
bottom
\end{verbatim}

When the callback returns, its JVM method frame disappears. The event loop then invokes the next callback. There is no parallel execution here. The state of pending work lives in heap objects inside the queue.

In a real Java event-loop framework, the queue is fed by socket readiness, timers, completed operations, or other system events. The principle remains the same: one event-loop thread executes one callback at a time.

\paragraph{C: \texttt{select}/\texttt{poll}/\texttt{epoll}-style event loops.}

C can implement concurrency without parallelism through non-blocking I/O and an event loop. The program stores state for each connection and uses the operating system to wait until one or more file descriptors become ready.

A simplified \texttt{select()}-style sketch:

\begin{codebox}[]{c}{C: select/poll/epoll-style event loops.}{sus05}
#include <sys/select.h>
#include <unistd.h>
#include <stdio.h>

void handle_stdin(void) {
    char buffer[256];
    ssize_t n = read(STDIN_FILENO, buffer, sizeof(buffer) - 1);

    if (n > 0) {
        buffer[n] = '\0';
        printf("input: %s", buffer);
    }
}

int main(void) {
    while (1) {
        fd_set read_set;

        FD_ZERO(&read_set);
        FD_SET(STDIN_FILENO, &read_set);

        int max_fd = STDIN_FILENO;

        select(max_fd + 1, &read_set, NULL, NULL, NULL);

        if (FD_ISSET(STDIN_FILENO, &read_set)) {
            handle_stdin();
        }
    }

    return 0;
}
\end{codebox}


The useful application is single-threaded servers, terminal programs, network daemons, embedded event loops, and systems that must manage many waiting I/O sources without one thread per source.

The stack behavior is ordinary C:

\begin{verbatim}
main()
    calls select()
        kernel blocks process until fd is ready
    select returns

main()
    calls handle_stdin()
        handle_stdin frame is created
        reads available data
        returns
    handle_stdin frame disappears

main()
    loops again
\end{verbatim}

There is no suspended C stack frame for each waiting file descriptor. Waiting state is stored in kernel data structures and in user-defined C structs.

A more realistic C event loop uses callbacks and explicit state:

\begin{verbatim}
struct watcher {
    int fd;
    void (*on_readable)(int fd, void *data);
    void *data;
};
\end{verbatim}

The callback supplies behavior; the \texttt{data} pointer supplies state. This is the C equivalent of a runtime task object:

\begin{verbatim}
watcher
+-----------------------------+
| fd                          |
| function pointer            |
| user data pointer           |
+-----------------------------+
\end{verbatim}

The event loop does not preserve function frames. It repeatedly waits for readiness, then calls the appropriate handler.

Duff's Device is not an event loop and not a concurrency mechanism by itself. It is useful only as a reminder that C can manually encode unusual control-flow/state-machine behavior. In C event loops, the usual serious technique is explicit state structs plus callbacks, not preserved stack frames.

\paragraph{Bash: limited shell-level interleaving and process-based contrast.}

Bash does not provide coroutine objects or a general single-threaded event loop like \texttt{asyncio}. It can still perform simple event-like polling in the shell itself.

Example:

\begin{verbatim}
while true; do
    if read -t 1 -r line; then
        echo "input: $line"
    else
        echo "tick"
    fi
done
\end{verbatim}

The useful application is simple supervisors, terminal scripts, polling loops, and small automation tasks.

The stack behavior is simple:

\begin{verbatim}
shell loop
    read with timeout
    if input exists, handle it
    otherwise run timeout branch
    loop again
\end{verbatim}

There is no saved coroutine frame. State lives in shell variables.

Bash pipelines and background jobs are different:

\begin{verbatim}
producer | consumer
\end{verbatim}

or:

\begin{verbatim}
long_task &
other_task &
wait
\end{verbatim}

These use operating-system processes. They may run in parallel if the OS schedules them on different cores. That is not single-threaded concurrency inside Bash. It is process-level concurrency managed by the operating system.

\paragraph{The core implementation idea.}

Concurrency without parallelism works by separating \textbf{waiting} from \textbf{executing}. A task that waits must leave behind a resumable description of what should happen later.

That description may be:

\begin{itemize}
    \item a Python coroutine frame stored in an \texttt{asyncio.Task};
    \item a JavaScript callback or promise continuation stored by the host/runtime;
    \item a PHP fiber object;
    \item a Java \texttt{Runnable} object in an event-loop queue;
    \item a C struct containing file descriptor, callback, and user data;
    \item a shell variable/loop state in Bash.
\end{itemize}

The single-threaded event-loop pattern is:

\begin{verbatim}
ready queue / event source
        |
        v
event loop chooses one ready task
        |
        v
task runs on the only thread
        |
        +--> finishes
        |
        +--> registers waiting operation and suspends/returns
        |
        v
event loop chooses next ready task
\end{verbatim}

Only one task has an active call stack at a time. Suspended or waiting tasks are represented by heap/runtime/kernel state, not by several simultaneously executing stacks on the same thread.

The comparison is:

\begin{center}
\begin{tabular}{|p{0.22\textwidth}|p{0.30\textwidth}|p{0.34\textwidth}|}
\hline
\textbf{Model} & \textbf{What is active?} & \textbf{Where waiting state lives} \\
\hline
Python \texttt{asyncio} & One coroutine executes Python bytecode at a time & Task/coroutine objects, Futures, event loop timers/I/O registrations \\
\hline
JavaScript / V8 & One JS call stack executes at a time & Host queues, promise continuations, callback references \\
\hline
PHP Fibers & One fiber runs at a time on the thread & Fiber objects preserve PHP execution context \\
\hline
Java event loop & One callback runs at a time & Heap queue of \texttt{Runnable} objects and framework event state \\
\hline
C event loop & One handler runs at a time & Kernel readiness state plus user structs and callback pointers \\
\hline
Bash polling loop & One shell command sequence runs at a time & Shell variables; external jobs are OS processes, not shell coroutines \\
\hline
\end{tabular}
\end{center}

The major limitation is CPU-bound work. If a task performs a long computation without yielding or returning to the event loop, it blocks the whole single-threaded system.

For example, in Python:

\begin{verbatim}
async def bad():
    while True:
        pass
\end{verbatim}

In JavaScript:

\begin{verbatim}
while (true) {
    // blocks the event loop
}
\end{verbatim}

In both cases, no other task can run on that thread. The operating system may still preempt the whole process and run other processes, but the runtime's internal tasks are stuck because the event loop has not regained control.

Therefore, concurrency without parallelism is excellent for waiting-heavy programs, not for CPU-heavy programs. It lets one thread manage many in-flight operations by storing suspended state and resuming tasks only when they can make progress.


\subsubsection{Asynchronicity: Non-Blocking Event-Driven I/O Execution Profiles, Engine Event Loops, and Single-Threaded Concurrency Subsystems}

\textbf{Asynchronicity} means that an operation is started now, but its result is handled later. The calling code does not wait by occupying the execution thread until the operation finishes. Instead, the program registers interest in the future result, returns control to a scheduler or event loop, and resumes the appropriate continuation when the result becomes available.

This is most useful for I/O. Disk, network, timers, user input, database requests, and subprocess output are slow compared with CPU execution. If a program blocks the only thread while waiting for such operations, no other task inside that thread can advance. Asynchronous execution avoids that by separating \emph{starting} an operation from \emph{handling} its completion.

The basic pattern is:

\begin{verbatim}
start operation
    operation is not ready yet
    register callback / future / task / continuation
    return to event loop

event loop waits for readiness or completion

operation becomes ready
    event loop resumes the waiting task
    result is processed
\end{verbatim}

Asynchronicity is therefore not the same as parallelism. A single-threaded asynchronous runtime may manage many in-flight operations, but it still executes only one callback or coroutine at a time on that thread.

\paragraph{Python / CPython: \texttt{asyncio}, \texttt{await}, tasks, and futures.}

Python's \texttt{asyncio} library provides a standard event-loop model. An \texttt{async def} function creates a coroutine. The coroutine runs until it reaches an \texttt{await}. At that point, it suspends and gives control back to the event loop.

\begin{codebox}[]{python}{Python / CPython: asyncio, await, tasks, and futures.}{sus27}
import asyncio


async def fetch_like_task(name, delay):
    print(name, "started")

    await asyncio.sleep(delay)

    print(name, "finished")


async def main():
    await asyncio.gather(
        fetch_like_task("A", 1),
        fetch_like_task("B", 1),
        fetch_like_task("C", 1),
    )


asyncio.run(main())
\end{codebox}


The useful application is waiting-heavy work: HTTP clients, websocket servers, database clients, queue consumers, API crawlers, and many open network connections.

The important line is:

\begin{verbatim}
await asyncio.sleep(delay)
\end{verbatim}

This does not block the whole thread. It registers a timer with the event loop and suspends the current coroutine. The event loop can then run another coroutine.

The runtime state is roughly:

\begin{verbatim}
event loop
+----------------------------------+
| ready task queue                 |
| timer registrations              |
| I/O readiness registrations      |
+----------------------------------+

Task A
+----------------------------------+
| coroutine object                 |
| suspended frame at await         |
| local variables                  |
| awaited Future / timer           |
+----------------------------------+
\end{verbatim}

The stack behavior is:

\begin{verbatim}
event loop resumes Task A
    Python frame for fetch_like_task runs
    reaches await
    frame is suspended into coroutine/task state
    control returns to event loop

event loop resumes Task B
    same pattern

later timer for Task A is ready
    event loop resumes Task A after await
\end{verbatim}

The suspended coroutine is not an active native stack frame while waiting. CPython preserves the Python-level frame state in runtime-managed coroutine/task structures. The native thread is free to return to the event loop.

A wrong version is:

\begin{verbatim}
import time

async def bad():
    time.sleep(5)
\end{verbatim}

This blocks the event-loop thread. Other coroutines cannot run during the sleep. In asynchronous code, waiting operations must be expressed through awaitable non-blocking operations.

\paragraph{JavaScript / V8: host event loops, AJAX, promises, and \texttt{async}/\texttt{await}.}

JavaScript asynchronous execution is normally provided by a host environment around the engine. V8 executes JavaScript, but the browser, Node.js, or Electron-style host provides timers, network operations, file-system APIs, event queues, and the event loop.

A browser example using \texttt{fetch}:

\begin{codebox}[]{javascript}{JavaScript / V8: host event loops, AJAX, promises, and async/await.}{sus16}
async function loadData() {
    console.log("before request");

    const response = await fetch("/data.json");
    const data = await response.json();

    console.log(data);
}

loadData();

console.log("after scheduling");
\end{codebox}


The useful application is AJAX-style loading: retrieving data from the server without freezing the page or reloading the document.

The control behavior is:

\begin{verbatim}
loadData starts
    logs "before request"
    starts fetch request
    reaches await
    loadData is suspended

top-level script continues
    logs "after scheduling"

network response arrives later
    host queues continuation
    event loop resumes loadData
\end{verbatim}

During the wait, the JavaScript call stack is empty or running other work. The old stack is not preserved as an active stack. The continuation is represented by promise and async-function runtime state.

Conceptually:

\begin{verbatim}
async function state
+----------------------------------+
| suspended instruction position   |
| local bindings                   |
| pending Promise                  |
+----------------------------------+
\end{verbatim}

A Node.js example:

\begin{codebox}[]{javascript}{JavaScript / V8: host event loops, AJAX, promises, and async/await.}{sus15}
const fs = require("fs/promises");

async function readFile() {
    console.log("before read");

    const text = await fs.readFile("notes.txt", "utf8");

    console.log(text);
}

readFile();

console.log("after call");
\end{codebox}


The file operation is started through Node's host APIs. JavaScript execution returns to the event loop. When the operation completes, the promise is resolved and the async function continues.

The event-loop side is:

\begin{codebox}[]{javascript}{JavaScript / V8: host event loops, AJAX, promises, and async/await.}{sus14}
host starts I/O operation
    stores promise continuation

I/O completes
    host queues microtask / continuation

V8 executes continuation
    async function resumes after await
\end{codebox}


Again, V8 executes the JavaScript frame only when resumed. The host provides the asynchronous I/O machinery.

\paragraph{PHP: request-oriented execution, Fibers, and async frameworks.}

Traditional PHP web execution is often synchronous: a PHP-FPM worker receives a request, runs the PHP script, blocks during ordinary database or file operations, returns a response, and then handles another request. However, modern PHP also supports coroutine-like async systems through Fibers and async frameworks.

A simple Fiber sketch:

\begin{codebox}[]{php}{PHP: request-oriented execution, Fibers, and async frameworks.}{sus18}
<?php

$fiber = new Fiber(function () {
    echo "before wait\n";

    Fiber::suspend("waiting");

    echo "after wait\n";
});

$value = $fiber->start();
echo $value . PHP_EOL;

$fiber->resume();
\end{codebox}


This is not yet real I/O, but it shows the mechanism: the PHP execution context can suspend and later resume.

For actual async frameworks, the pattern is:

\begin{verbatim}
start non-blocking socket / timer / HTTP request
    suspend current fiber or promise continuation
    return to event loop

I/O ready or complete
    event loop resumes fiber / continuation
\end{verbatim}

At the Zend Engine level, a Fiber preserves a PHP execution context:

\begin{verbatim}
Fiber object
+----------------------------------+
| suspended PHP call stack         |
| local variables                  |
| current opcode position          |
| status                           |
+----------------------------------+
\end{verbatim}

The useful applications are high-concurrency servers, websocket systems, async HTTP clients, message consumers, and event-driven workers. The key difference from ordinary PHP request execution is that one worker process can keep many operations in flight and resume them when ready.

\paragraph{Java: NIO selectors, event loops, and completion callbacks.}

Java has several asynchronous styles. One low-level model is Java NIO, where a selector waits for readiness on multiple channels. A single thread can watch many sockets and process whichever ones become ready.

A simplified shape is:

\begin{verbatim}
Selector selector = Selector.open();

while (true) {
    selector.select();

    for (SelectionKey key : selector.selectedKeys()) {
        if (key.isReadable()) {
            // read available data
        }

        if (key.isWritable()) {
            // write pending data
        }
    }

    selector.selectedKeys().clear();
}
\end{verbatim}

The useful application is scalable network servers. Instead of one blocking thread per connection, the program keeps connection state in heap objects and lets the selector report which channels can make progress.

The stack behavior is:

\begin{verbatim}
event-loop thread
    calls selector.select()
        blocks in kernel until at least one channel is ready
    returns to Java code
    invokes handler logic for ready channels
    handler returns
    loop repeats
\end{verbatim}

There is no preserved Java method frame per socket. The per-connection state lives in objects:

\begin{verbatim}
ConnectionState object
+----------------------------------+
| socket/channel reference         |
| input buffer                     |
| output buffer                    |
| parser state                     |
+----------------------------------+
\end{verbatim}

Java also supports future/completion-style async code:

\begin{verbatim}
CompletableFuture
    .supplyAsync(() -> loadData())
    .thenApply(data -> transform(data))
    .thenAccept(result -> System.out.println(result));
\end{verbatim}

This is useful for composing operations that complete later. However, this example may use worker threads from a pool, so it may involve parallelism depending on the executor. It is asynchronous, but not necessarily single-threaded. The key architectural idea is still delayed continuation: the next function is invoked when the previous result becomes available.

\paragraph{C: non-blocking file descriptors and explicit event loops.}

C has no built-in async/await syntax. Asynchronous I/O is usually built with operating-system APIs and explicit state machines.

A simplified non-blocking event loop uses readiness APIs such as \texttt{select}, \texttt{poll}, \texttt{epoll}, \texttt{kqueue}, or platform-specific completion mechanisms.

Conceptual C structure:

\begin{codebox}[]{c}{C: non-blocking file descriptors and explicit event loops.}{sus04}
struct connection {
    int fd;
    char input_buffer[4096];
    size_t input_used;
    void (*on_readable)(struct connection *);
};

while (1) {
    wait_for_ready_file_descriptors();

    for each ready connection:
        connection->on_readable(connection);
}
\end{codebox}


The useful application is high-performance servers, proxies, terminal multiplexers, embedded event systems, and network daemons.

The stack behavior is ordinary C:

\begin{verbatim}
main event loop
    waits in kernel readiness function
    returns with ready fd list
    calls on_readable(connection)
        handler frame is created
        reads available bytes
        updates connection struct
        returns
    handler frame disappears
    loop continues
\end{verbatim}

There is no suspended C function frame for each connection. The state must be stored explicitly:

\begin{verbatim}
connection struct
+----------------------------------+
| file descriptor                  |
| buffers                          |
| parser state                     |
| callback pointer                 |
| user data                        |
+----------------------------------+
\end{verbatim}

Duff's Device and similar tricks can be used to hand-write state-machine-like code, but serious C asynchronous systems usually rely on explicit state structs, callbacks, readiness APIs, or non-portable coroutine/fiber libraries.

A key distinction in C systems is readiness versus completion:

\begin{itemize}
    \item \textbf{readiness}: the OS says a descriptor can be read or written without blocking;
    \item \textbf{completion}: the OS says an operation that was submitted earlier has completed.
\end{itemize}

APIs such as \texttt{select}, \texttt{poll}, and \texttt{epoll} are readiness-oriented. Some Windows I/O models and newer systems can be completion-oriented. The event-loop design changes, but the stack rule remains similar: waiting state is stored outside ordinary call frames.

\paragraph{Bash: asynchronous-looking process orchestration, not runtime async.}

Bash does not provide an async/await runtime. It can run background processes and wait for them later:

\begin{verbatim}
long_task_one &
pid_one=$!

long_task_two &
pid_two=$!

wait "$pid_one"
wait "$pid_two"
\end{verbatim}

The useful application is shell-level orchestration: start several external commands, let the operating system run them, and collect their exit statuses.

This is not single-threaded runtime asynchronicity. Bash creates child processes. Each child has its own process state, native stack, and scheduler identity. The operating system may run them in parallel on multiple cores.

Bash can also use polling:

\begin{verbatim}
while true; do
    if read -t 1 -r line; then
        echo "input: $line"
    else
        echo "tick"
    fi
done
\end{verbatim}

This is useful for small interactive scripts, but it is not a general asynchronous runtime. The shell loop remains explicit. State lives in shell variables and child process status, not in coroutine objects.

\paragraph{Event-loop architecture.}

A single-threaded asynchronous event loop has four main responsibilities:

\begin{itemize}
    \item track operations that are waiting;
    \item ask the operating system or host which operations are ready or complete;
    \item place ready continuations into a queue;
    \item execute one continuation at a time.
\end{itemize}

The structure is:

\begin{verbatim}
event loop
+----------------------------------+
| ready queue                      |
| timers                           |
| I/O registrations                |
| pending futures/promises/tasks   |
+----------------------------------+
        |
        v
choose one ready continuation
        |
        v
run it until it returns or suspends
        |
        v
repeat
\end{verbatim}

Only one callback or coroutine owns the thread at a time. If it returns quickly or awaits properly, the system remains responsive. If it runs CPU-heavy code for too long, the event loop is blocked.

\paragraph{Comparison.}

\begin{center}
\begin{tabular}{|p{0.18\textwidth}|p{0.32\textwidth}|p{0.34\textwidth}|}
\hline
\textbf{Language/runtime} & \textbf{Async mechanism} & \textbf{Stack and state model} \\
\hline
Python / \texttt{asyncio} & Coroutines, Tasks, Futures, selector/proactor event loop & Coroutine frames suspend into Task/Future state; one event-loop thread executes one coroutine at a time \\
\hline
JavaScript / V8 & Promises, \texttt{async}/\texttt{await}, host event loop, callbacks & Async function state and promise continuations are stored by engine/host; callback stack is created later \\
\hline
PHP / Zend & Fibers and async frameworks & Fiber preserves PHP execution context; event loop resumes fibers when I/O is ready \\
\hline
Java & NIO selectors, event-loop frameworks, futures/completion stages & Event-loop handlers use normal JVM frames; connection state lives in heap objects; futures may use thread pools \\
\hline
C & Non-blocking descriptors, callbacks, explicit state machines, OS readiness/completion APIs & No suspended C frame; state stored in structs; handlers create ordinary native frames when called \\
\hline
Bash & Background jobs, \texttt{wait}, polling loops & Mostly process orchestration; child processes have separate stacks; shell has no coroutine runtime \\
\hline
\end{tabular}
\end{center}

The central architectural rule is:

\begin{quote}
Asynchronous systems do not make waiting disappear. They move waiting out of the active call stack and into runtime, host, kernel, or heap-managed state.
\end{quote}

That is the real mechanism. A blocking program waits by keeping the current call path occupied. An asynchronous program records what should happen later, returns to the event loop, and lets other ready work run. When the external event completes, the runtime reconstructs or resumes the appropriate continuation.

This is why async programming belongs in the same family as generators and coroutines. All of them depend on runtime-managed continuation state. The difference is purpose: generators suspend to produce values lazily; coroutines suspend to cooperate with other execution units; asynchronous tasks suspend because the next step depends on an external event that is not ready yet.